\chapter{Related Work 相關研究}

\section{AI 技術相關研究}

\subsection{自然語言需求與可驗證規劃}

TravelPlanner 以 1,225 筆規劃意圖與參考方案評估語言代理，指出多重條件、資訊蒐集與長流程規劃的困難。後續研究將需求轉為形式化限制，再交由求解工具處理，在其 TravelPlanner 設定下達到 93.9\% 成功率。這些結果支持將語言理解與限制驗證分工，但不同論文的模型、工具與評估設定必須分開閱讀，不能將他處成功率當成本研究的預期成果~\cite{llm_planning,llm_planning_tools}。

TTG 也採取自然語言轉符號表示、再由混合整數線性規劃求解的架構。另一方面，HyperTree Planning 以階層式規劃結構分解任務並反覆修訂，顯示模型端的規劃方式仍可改進。因此，本研究選擇程式化查核，是基於條件需可追溯、規則需一致的設計需求，並不推論所有語言模型規劃皆無效~\cite{ttg,hypertree_planning}。

\subsection{語意檢索與結構化證據}

BRIGHT 將需要推理的檢索需求獨立評測，提醒本研究不能假設語意向量在所有查詢皆優於關鍵字。混合檢索研究則以稀疏與稠密方法的互補性改善檢索。本研究因此比較 BM25、向量與混合召回，讓「適合聚餐」等語意偏好與地名、料理名等明確詞彙皆有相應的檢索方法~\cite{bright,hybrid_retrieval}。

K-RagRec 使用知識圖譜提供推薦所需的結構化資訊，並在電影與書籍任務中評估推薦表現。本研究參考其以結構化證據支持推薦的設計，將圖譜關係設計為餐廳、屬性、來源片段與有效時間；是否需要圖譜，則以相同事實與規則的關聯表版本比較。保存證據與使用圖譜是兩個不同因素，實驗須分別控制~\cite{knowledge_graph_llm_recommendation}。

\section{餐廳與公車的領域知識相關研究}

\subsection{主觀偏好與用餐情境}

餐廳的氛圍、適合對象與用餐目的具有主觀性。主觀屬性研究指出，同一句描述可能依使用者觀點產生不同解讀；餐廳對話研究則將使用者批評轉為正向偏好，再檢索相關評論支持推薦。這些成果支持保留「聚餐」「安靜」等原始語句與來源片段，並容許使用者修正偏好，而不把氛圍標籤當成客觀保證~\cite{subjective_attributes,critique_preference}。

在位置與時間方面，Habib 等人將目前位置、請求時間、簽到歷史及人氣整合為分數，並依簽到頻率推估餐廳的活動時段。該研究使用 Foursquare 的巴黎餐廳子集，以兩個查詢案例展示推薦結果；作者在結論中明確指出尚未納入顧客評論，實驗段落亦未報告整體推薦準確率或統計檢定。STiSAN 則以時空間隔建模下一個地點的選擇。這些方法說明位置與時間會影響推薦，卻不能使「偏好相似」直接等同於「到店後能完成用餐」。本研究將這兩種問題分開：偏好由相關性排序處理，營業與交通則由可行性查核處理~\cite{restaurant_habib,poi_stisan}。

\subsection{公車時間、轉乘與動態重規劃}

公車可行性須依時間順序處理。到達上車站之前已離站的車輛，不能因為 API 回傳正常狀態而視為可搭乘；轉乘則須先到達前段下車站，完成站間步行與緩衝，才能搭上後段班次。動態重規劃研究以班次、連接段與步行關係描述這些限制，並在 Perth、Berlin、Paris、London 四個運輸網路比較使用者主動請求與系統主動更新的方法。本研究以其中的步行、候車與轉乘模型分析候選餐廳的交通便利性；旅途中重規劃留作後續擴充~\cite{bus_dynamic_replanning}。

此外，公車到站時間是預測值。以 Stockholm 公車資料進行的研究，透過 Student-t 模型與貝氏推論表示延誤的不確定性，說明評估不應只關注單一時間預估。本研究不另訓練到站模型，而是保存外部預估的時間依據、加入可調整的緩衝，並將「查核通過」限定為在所採資料與規則下成立~\cite{bus_delay_uncertainty}。

TDX 的旅運規劃與公車介面，分別提供路線候選及公車資料。系統實作時須核對站位、方向、子路線、班次可識別性、時間單位與更新欄位，避免只憑路線名稱連結資料。旅運規劃回傳的候選是否有當下可搭乘班次，仍交由後續查核處理~\cite{tdx_planner,tdx_bus}。

\section{Article Summary Table}

表~\ref{table:article_summary} 將 AI 技術與餐廳、公車領域研究各列六篇。


{\small\singlespacing
\renewcommand{\arraystretch}{1.25}
\setlength{\tabcolsep}{3pt}
\begin{longtable}{|P{\dimexpr 0.0650\linewidth-2\tabcolsep-1.1250\arrayrulewidth\relax}|P{\dimexpr 0.0650\linewidth-2\tabcolsep-1.1250\arrayrulewidth\relax}|P{\dimexpr 0.1300\linewidth-2\tabcolsep-1.1250\arrayrulewidth\relax}|P{\dimexpr 0.1300\linewidth-2\tabcolsep-1.1250\arrayrulewidth\relax}|P{\dimexpr 0.1350\linewidth-2\tabcolsep-1.1250\arrayrulewidth\relax}|P{\dimexpr 0.1500\linewidth-2\tabcolsep-1.1250\arrayrulewidth\relax}|P{\dimexpr 0.1600\linewidth-2\tabcolsep-1.1250\arrayrulewidth\relax}|P{\dimexpr 0.1650\linewidth-2\tabcolsep-1.1250\arrayrulewidth\relax}|}
\caption{Article Summary Table}\label{table:article_summary}\\
\hline
\textbf{Ref.} & \textbf{Year} & \textbf{Problem} & \textbf{AI Method} & \textbf{Domain}\newline\textbf{Knowledge} & \textbf{Dataset} & \textbf{Findings} & \textbf{Limitation} \\ \hline
\endfirsthead
\hline
\textbf{Ref.} & \textbf{Year} & \textbf{Problem} & \textbf{AI Method} & \textbf{Domain}\newline\textbf{Knowledge} & \textbf{Dataset} & \textbf{Findings} & \textbf{Limitation} \\ \hline
\endhead
% Limitation evidence: ICML final PDF p.9, Impact Statement: commonsense criteria based on authors consensus.
\cite{llm_planning} & 2024 & 多重限制旅遊規劃 & 語言代理與工具使用評測 & 旅遊安排與限制 & TravelPlanner，1,225 筆規劃意圖 & GPT-4 在其兩階段設定的最終通過率為 0.6\% & 常識性評估規則依作者共識制定，未涵蓋不同使用者的全部常識判準 \\ \hline
% Limitation evidence: Sec.7 Limitation, p.3442: Prompt Designing, Solver Runtime, Risky Data.
\cite{llm_planning_tools} & 2025 & 可驗證的限制規劃 & LLM 轉形式化限制，再由求解器處理 & 旅遊條件與不可滿足條件診斷 & TravelPlanner、作者建立的國際旅遊資料 & 其設定下成功率達 93.9\% & 提示與編碼設計耗時；求解可能逾時，且無法辨識資料庫中的錯誤或不安全資訊 \\ \hline
% Limitation evidence: Sec.5.3 and Sec.6, pp.244-246: subjective proxy evaluation, correct parsing prerequisite, personalization and multi-round dialog.
\cite{ttg} & 2024 & 自然語言旅遊需求求解 & 微調 LLM、混合整數線性規劃 & 航班、住宿與景點 & 依真實統計建立的合成資料 & 合成測試的回譯精確比對約 91\%；另有使用者評估 & 可行性與最適性以需求正確解析為前提；個人化與多輪對話仍待研究 \\ \hline
% Limitation evidence: Appendix G, p.25: subjective relevance judgments and multimodal settings.
\cite{bright} & 2025 & 需要推理的檢索 & BM25、向量等方法評測 & 多領域推理需求 & BRIGHT，1,384 筆查詢 & 顯示推理型檢索仍具挑戰 & 相關性標註具主觀性；尚未處理多模態檢索 \\ \hline
% Limitation evidence: Sec.6 Limitations, p.13365: sparse relevance labels, false negatives and exact-match bias.
\cite{hybrid_retrieval} & 2023 & 稀疏與稠密檢索互補 & 混合檢索訓練目標 & 通用資訊檢索 & MS MARCO、Natural Questions、TREC Robust04 & 互補性目標改善其檢索評測 & 稀疏相關性標註可能造成偽負例，且 MS MARCO 標註偏向精確詞彙匹配 \\ \hline
% Limitation evidence: Sec.6 Limitations, p.27160: 7B/8B, Freebase, popularity-based retrieval policy.
\cite{knowledge_graph_llm_recommendation} & 2025 & LLM 推薦的知識支持 & K-RagRec & 電影與書籍推薦 & MovieLens、Amazon Book & 結構化知識改善推薦；另分析虛構項目推薦 & 僅評估 7B／8B 模型及 Freebase；選擇性檢索策略以熱門度為依據 \\ \hline
% Limitation evidence: Introduction and Conclusion, pp.12287,12292: preliminary sketch, formal semantics and evaluation methodology left open.
\cite{subjective_attributes} & 2022 & 主觀屬性的表達 & 概念框架 & 推薦語句的觀點與主觀性 & 無；概念性論述 & 整理主觀屬性問題與研究方向 & 僅提出初步框架；完整形式語意、方法與評估指標仍屬待研究項目 \\ \hline
% Limitation evidence: Sec.4, p.1918: exclusion and clarification of critique scope are future work.
\cite{critique_preference} & 2021 & 餐廳批評如何轉為偏好 & 語言模型轉換、嵌入與評論排序 & 用餐氛圍與偏好修正 & 餐廳批評與評論資料 & 消融顯示批評轉換可改善其推薦 & 尚未處理以批評排除不適合餐廳，以及批評範圍歧義的澄清策略 \\ \hline
% Limitation evidence: full paper Sec.III.A.4 p.318, Sec.IV pp.318-319 (two illustrated requests; no aggregate accuracy/significance test reported), Sec.V p.320 (authors explicitly exclude customer reviews).
\cite{restaurant_habib} & 2016 & 位置、時間與偏好的餐廳推薦 & 時間衰減、互相強化與乘積評分 & 地理距離、簽到推估活動時段與人氣 & Foursquare 巴黎子集：7,991 間餐廳、111,325 筆簽到、6,903 名使用者 & 兩個查詢案例呈現偏好與時段的作用；早餐案例較比較方法符合用餐時段 & 未納入顧客評論；評估僅展示查詢案例，未報整體準確率或統計檢定 \\ \hline
% Limitation evidence: Sec.IV.B.2 Finding 5, p.2094: TAAD sometimes useful; omitted last-POI/candidate spatial intervals offered as a possible reason.
\cite{poi_stisan} & 2022 & 下一個 POI 的推薦 & STiSAN 時空注意力 & 到訪的時空間隔 & Gowalla、Brightkite、Weeplaces、Changchun & 評估時空間隔建模對序列推薦的作用 & TAAD 僅在部分情況有效；作者推測與忽略序列末項和候選 POI 的空間間隔有關 \\ \hline
% Limitation evidence: Sec.5 Delay Modelling and Sec.6, pp.8435-8436: synthetic delays; real-world delay data and multicriteria routing as future work.
\cite{bus_dynamic_replanning} & 2025 & 公共運輸中的動態重規劃 & 連接掃描、相關時刻表範圍與 push／pull 更新 & 班次、轉乘、步行與延誤 & Perth、Berlin、Paris、London 的 GTFS 網路與合成延誤 & 在其模擬設定下，改善抵達時間並降低更新查詢負擔 & 延誤採合成模型；真實延誤驗證與多目標路由仍待研究 \\ \hline
% Limitation evidence: Sec.V.D and VI: overparameterization at end-of-route stops, broader datasets and network/high-low frequency comparisons needed.
\cite{bus_delay_uncertainty} & 2022 & 公車延誤與預測不確定性 & Student-t 模型、貝氏推論 & 候車可靠性、時段與站間關係 & Stockholm 高頻公車資料 & 以機率預測表達延誤；其 Student-t 模型優於比較的 Gaussian 模型 & 需更大範圍資料驗證網路效應與高／低頻差異；末段站位可能有過度參數化 \\ \hline
\end{longtable}
}


\section{現有研究、限制與研究缺口}

從原研究的限制可歸納三項設計課題。第一，形式化求解仍仰賴問題編碼及資料正確性；第二，檢索標註具有主觀性與不完整性，相關性分數不能被解讀為客觀事實的保證；第三，動態交通方法在合成延誤下的成效，仍需透過真實延誤資料驗證。這些是原研究明示的限制，與是否涵蓋臺中餐廳場域為不同層次的問題~\cite{llm_planning_tools,bright,hybrid_retrieval,bus_dynamic_replanning}。

現有研究已提供語意檢索、主觀偏好、時空地點推薦與公共運輸規劃等方法。餐廳情境推薦關心語意是否相符，交通規劃則分析如何抵達；把兩者接起來，還需要決定交通條件如何影響候選餐廳的取捨~\cite{critique_preference,bus_dynamic_replanning}。

本研究關注的缺口是：在約會、慶生等用餐情境中，先從餐廳評論召回語意相符的候選，再依臺中公車的步行、候車、乘車與轉乘條件篩選餐廳，是否能比只看語意分數提供更合適的最終推薦？此外，若交通或店家證據不足，如何避免把暫時無法判斷的候選當成已確認可行？

上述文獻尚未直接回答本研究所界定的臺中場域問題，因此本研究以情境候選集合、交通查核與餐廳排序之間的接合方式作為設計重點。研究定位為推薦流程與場域驗證，不以未找到完全相同的題目作為首次提出或既有研究不存在的證明。
